\begin{lemma}[Restoring the terminal hierarchy from two certificates]
Fix a counterexample $C$ on the exact tree $T$, its minimum context,
and a raw Case Three datum $D$. Let $\beta$ be missing on
$B=T(y_{p-2})$ and not missing at $y_p$. Put
$B^+=T(y_{p-1})$, using the full-tail prefix definitions.
Let $\sigma$ be the $(D.\alpha,\beta)$-component swap rooted at $y_p$.
Suppose $F$ is \noderef{P2752_T2_PrefixData} on $B^+$ for $\sigma$,
every displayed edge of $T$ is fixed, and $\mathcal C^T$ is any
full-tree corresponding certificate from the original ETT to $\sigma$
on $T$. Then $\mathcal C^T$ admits a hierarchy base and a good hierarchy
on $T$ which satisfy SameGoodHierarchy with $C.\mathit{good}$.
No equality between $\mathcal C^T$ and the prefix certificate is assumed.
\end{lemma}
\begin{proof}
Write $H=C.\mathit{hierarchy}$ and $q=H.\mathit{dividerCount}$.
The notation $B,B^+$ is given by \noderef{HierarchyFullTailPrefix} and
\noderef{HierarchyFullTailLength}. The inequality $\ell\ge2$ in
\noderef{HierarchySituationTwoCommonTailData} gives
$|B|=|T|-2$ and $|B^+|=|T|-1$; these are exact prefixes of $T$ by
\noderef{TreeSequencePrefix}. The fields of $F$ from
\noderef{P2752_T2_PrefixData} will be denoted as follows: the old prefix
ETT, base, hierarchy and good witness are $\mathcal E^+,b^+,H^+,g^+$;
the new prefix corresponding certificate, base, hierarchy and good
witness are $\mathcal C^+,b^+_\sigma,K^+,g^+_\sigma$.
The old prefix ETT has the original rung and the identical certified
series. Both corresponding certificates in
\noderef{CorrespondingExtendedTashkinovTree} therefore have the original
positive ladder trees, raw stored sets, connecting-edge sets, and kinds.
This identifies their terminal ladder trees with the old $T_n$ and their
terminal kinds with the old kind. It does not identify their certified
series or internal witnesses.

We first check the only witness equality required for a common base.
Suppose the terminal kind is PE. Inspect the old, new prefix and new full
terminal certificates in \noderef{TashkinovSeries}. In each case
\noderef{TashkinovSeriesStep} forces the PE constructor: an initial
certificate has kind none and SE and RE have their respective different
kinds. Write $(v,f,u)$ for the vertex, exact edge and outside endpoint of
the old \noderef{MaximumDefectiveChoice}, and $(v',f',u')$ for either
replacement choice. The constructor's singleton-edge equations and
connecting-edge agreement give $\{f'\}=\{f\}$, whence $f'=f$.
The choice fields give
\[
 G[f]=\{v,u\}=\{v',u'\},\qquad
 v,v'\in V(T_n),\quad u,u'\notin V(T_n).
\]
Equality of these unordered endpoint pairs has two possibilities. In the
uncrossed one $v'=v$. In the crossed one $v'=u$, contradicting the inside
and outside memberships just displayed. Thus $v'=v$ in both replacement
certificates. The chosen-vertex equations of their PE constructors give
the same terminal chosen vertex for the old, prefix and full certificates.
This conclusion used exact singleton identities and endpoint splits, not
an unwarranted equality of arbitrary corresponding witnesses.

Construct a base $b^T_\sigma$ for $\mathcal C^T$ using
\noderef{HierarchyBase}. Copy the ordered core and closure of
$b^+_\sigma$ and transport its support position across equality of the
terminal ladder trees, retaining its numerical value. Support positivity
copies. The terminal missing-color equation is reflexive, since both the
base coloring and the replacement generating coloring are $\sigma$.
If the terminal kind is PE, the chosen-vertex equality just proved
transports the support-vertex clause; equality of the terminal ladder trees
and numerical support positions identifies their exact support prefixes
by \noderef{TreeSequencePrefix}. The prefix base's
\noderef{TashkinovClosure} and \noderef{OrderedTreeUnion} certificates
then give the full base's closure and union clauses. If the kind is not
PE, both replacements have that same non-PE kind, and the prefix base's
nonparallel clause identifies its copied core with $T_n$. This gives the
full base's nonparallel clause; its parallel implication is vacuous.
Thus every base field is supplied.

By \noderef{SameGoodHierarchy}, the prefix same-good relation identifies
$b^+_\sigma.\mathit{core}$ with $b^+.\mathit{core}$, and, in PE, their
closures. The old-base equalities in $F$ identify these with the original
core and closure. Hence the new full base has the original core and the
original closure whenever the terminal kind is PE.

Construct $K^T$ using \noderef{TashkinovHierarchy}. Its divider count is
that of $K^+$. Copy every nonterminal ordinary stage and replace only the
terminal stage $B^+$ by $T$. Copy each divider position by its numerical
value, now as an index of $T$; its bound follows from $B^+\preccurlyeq T$.
The base-stage equation follows from terminal-ladder equality. The terminal
stage equation is the definition. Below the terminal stage, the nested and
strict fields are those of $K^+$. For the last nesting compose its old
inclusion into $B^+$ with $B^+\preccurlyeq T$, comparing retained entries
as in \noderef{TreeSequenceInitialSegment}; its strict length inequality
follows from the old strict inequality and $|B^+|<|T|$. Divider positivity
and increasing order are numerical and copy. Each divider-prefix equation
copies because an exact prefix of $B^+$ at the same numerical index is
also that exact prefix of $T$. The core-inside-first field copies for
$q>0$; for $q=0$ its old target was $B^+$, so compose with its inclusion
in $T$. These checks give all hierarchy fields.

The prefix same-good relation identifies all ordinary stages, dividers
and cores of $K^+$ and $H^+$. The old comparisons in $F$ then identify
every nonterminal stage and numerical divider of $K^T$ with $H$, and the
terminal stages are both literally $T$. Counts, stage families and divider
families therefore agree exactly. Together with core, terminal-kind and
conditional closure agreement, this gives
\noderef{SameHierarchyStages} for $H,K^T$.

We check all six fields of \noderef{GoodHierarchy}. First identify the
computed color data. For every internal index, the ordinary stages agree
with the original stages and are contained in $B^+$. The definition
\noderef{HierarchyStarredStage}, using the common core and the retained
ordinary stages, gives the same assertion for starred stages. Vertex
sets here are the displayed sets of \noderef{TreeSequenceVertexSet}. The
vertex clause of \noderef{TreeColoringInvariant}, recorded in $F$,
identifies all their missing sets, including that on $T_n$.
Apply \noderef{P2752_T2_TerminalStoredColors} to both corresponding
certificates. Their old series are identical, so their terminal computed
\noderef{TashkinovStoredColors} equal the original computed set. This
uses terminal-ladder missing-set stability as well as equality of the
raw stored arrays. The subtraction defining
\noderef{HierarchyStoredColorsAt} therefore gives identical stored sets
on $H,K^T$, and \noderef{HierarchyStageMissingColors} gives identical
stage-missing sets. The same argument identifies the new full stored
sets with those of $K^+$, since their internal stages, base core and
terminal ladder are identical. Define the full Gamma function by copying
$g^+_\sigma.\mathit{gamma}$. For each relevant stored color, prefix
same-good and the old Gamma comparisons in $F$ identify this value with
$C.\mathit{good}.\mathit{gamma}$.

To check avoidance at the last stage we cannot merely restrict the prefix
claim. We prove equality of all full first-missing positions. For a color
other than $D.\alpha$ and $\beta$, the swap changes no incidence of that
color, so its missingness at every vertex of $T$ is unchanged, by the
pointwise formula of \noderef{TwoColorComponentSwap} and
\noderef{MissingColors}. For $D.\alpha$, the raw earlier missing witness
in \noderef{HierarchySituationTwoCaseThreeData} has index strictly before
the penultimate vertex, hence lies in $B$. For $\beta$, the assumed
tree-missing membership on $B$, expanded by \noderef{TreeMissingColors},
provides such a witness as well. For either affected color there is thus
a candidate missing position in $B\subseteq B^+$. Invariance preserves
all candidates in $B^+$. The least old and new candidates are at most
this preserved witness and therefore both lie in $B^+$; comparing the
candidate predicates there proves they are the same index. In the
unaffected case candidate sets agree everywhere, so their minima and their
empty-set fallbacks agree. This verifies the entire totalized definition
\noderef{HierarchyFirstMissingPosition}.

Every displayed edge of $T$ is fixed by hypothesis. Each edge of a copied
internal stage is that same exact edge of $T$, by its prefix relation.
The identical first-missing positions, ordinary and starred stages and
edge values give equality of all old and new full
\noderef{HierarchySegmentColors}, directly at each edge index in its
definition. For a relevant stored color, original goodness now gives
the full avoidance inclusion: the Gamma pair, stage-missing set and
segment-color set are all exactly those of the original inclusion.
This also checks the possible last-edge contribution explicitly.

The two-colors and disjoint-pairs fields copy from $g^+_\sigma$, because
the internal computed stored sets and Gamma function are identical.
The unique-refresh field copies, including uniqueness: all its quantified
stored sets, Gamma values, difference of the two stage vertex sets and
the coloring $\sigma$ are identical in the prefix and full presentations.
The parallel-base field copies because the full and prefix terminal
kinds agree and their terminal ladder, core, closure, index-zero stored
set and Gamma function agree. The positive-stage minus-closedness fields
copy because their exact stages, coloring, stored sets and Gamma control
sets agree. With the avoidance check above these are all six fields,
giving a full good witness $g^T_\sigma$.

Finally the same-stage relation already checked, equality of every internal
stored set, and equality of Gamma values at each member of that set are
exactly the two clauses of \noderef{SameGoodHierarchy}. They give
SameGoodHierarchy between $C.\mathit{good}$ and $g^T_\sigma$ and complete
the construction.
\end{proof}
